\clearpage
\begingroup
\fancyhead[L]{\small BEMERKUNGEN}
\renewcommand{\headrulewidth}{0pt}
\SourceLeaf{607}
\needspace{6\baselineskip}
\pdfbookmark[1]{Bemerkungen zu den Störungen der Pallas}{pallas-bemerkungen}
\begin{center}\large BEMERKUNGEN.\end{center}
\begingroup
\newcommand{\BIntroJ}{\text{\mdseries\AstroSymbols ♃}}
\newcommand{\BIntroP}{\text{\mdseries\AstroSymbols ⚴}}

\par\noindent\hspace*{\parindent}Von \textsc{Gauss}’ Untersuchungen über die Störungen der Pallas finden sich im Nachlass fast ausschließlich
numerische Rechnungen ohne erklärenden Text. Sie stehen auf einer großen Menge einzelner Zettel,
die im Archiv fünf Kapseln füllen. Die Berechnung der Oppositionen und den Vergleich mit den
Beobachtungen hat \textsc{Gauss} jedoch auch vielfach in seinen Handbüchern ausgeführt. Die große Menge des
Materials erschwerte die Sichtung erheblich. Schließlich gelang es jedoch, insbesondere durch das gleichzeitige
Studium des Briefwechsels, das Ganze in eine zusammenhängende Form zu bringen.

\par\noindent\hspace*{\parindent}Vorangestellt sind die Stellen aus dem Briefwechsel (S.~413—438), die den
historischen Verlauf der \textsc{Gauss}schen Untersuchungen deutlich machen. Über ihren Zusammenhang mit dem später
abgedruckten Text sei Folgendes bemerkt:

\medskip
\par\noindent\hspace*{\parindent}\textsc{Gauss} hatte während der Berechnung der Ceresstörungen den Plan gefasst, eine Tafel herzustellen,
die die Entwicklung der Störungsfunktion für alle bis dahin entdeckten kleinen Planeten
erleichtern sollte. Er schreibt darüber an \textsc{Olbers} am 3.~September sowie am 7.~und 29.~Oktober 1805. Auch findet
sich im Handbuch B\,e eine Reihe von Aufzeichnungen zur Methode, die \textsc{Gauss} hier anwenden wollte.
Die Tafel, bei deren Berechnung \textsc{Bessel} half, wurde fertiggestellt und ist im Nachlass vorhanden.
Da \textsc{Gauss} dann aber bei der Berechnung der Pallasstörungen andere Methoden anwandte,
kam sie nie zur Anwendung. Deshalb wurde auf ihren Abdruck verzichtet.

\par\noindent\hspace*{\parindent}Die Briefe an \textsc{Olbers} von Januar 1806 bis Oktober 1810 beziehen sich auf Vorarbeiten zur Berechnung
der Störungen, und zwar teils auf Untersuchungen darüber, wie genau sich die Beobachtungen
mit der Methode der kleinsten Quadrate darstellen lassen, teils auf die analytischen
Grundlagen für die Entwicklung der Störungsfunktion. Ende Oktober 1810 begann \textsc{Gauss} dann die erste Rechnung
zu den speziellen Störungen der Pallas (abgedruckt S.~473), auf die sich die Briefe vom
November 1810 beziehen. Der Brief an \textsc{Olbers} vom 13.~Dezember 1810 betrifft die zweite Rechnung der
speziellen Störungen (S.~483). Die Briefe von August und Oktober 1811 betreffen die erste Rechnung der allgemeinen
Störungen (S.~489). Im April 1812 begann \textsc{Gauss} die zweite Rechnung der allgemeinen Störungen,
die im Juli 1813 abgeschlossen wurde. Aus der Zeit von März 1812 bis Januar 1813 konnten neben dem Brief an
\textsc{Bessel} vom 5.~Mai 1812 nur die Briefe von \textsc{Olbers} an \textsc{Gauss} abgedruckt werden, da die entsprechenden Briefe
von \textsc{Gauss} an \textsc{Olbers} verloren gegangen sind. Dazu gehört auch der Brief,
in dem \textsc{Gauss} wie im genannten Brief an \textsc{Bessel} seine Entdeckung des rationalen Verhältnisses
der Umlaufzeiten von Pallas und Jupiter mitteilt. \textsc{Gauss} machte diese Entdeckung also erstmals vor der
Ausführung der zweiten Rechnung der allgemeinen Störungen im März oder Anfang April 1812 gemacht,
als er noch gar nicht mit dem kleinen Divisor \(18\BIntroJ-7\BIntroP\) gerechnet hatte. Es bleibt unklar (vgl. aber die
Ausführungen weiter unten auf S.~604), wie er das rationale Verhältnis in diesem Fall gefunden hat. Die Notizen,
\par\smallskip\noindent VII.\hfill 76\par
\noindent\SourceLeaf{608}die sich darüber im Nachlass fanden und auf S.~558 abgedruckt sind, stammen aus späterer Zeit, jedenfalls
\textls[120]{nach} der Ausführung von mindestens einem großen Teil der zweiten Rechnung.

\par\noindent\hspace*{\parindent}Besondere Sorgfalt verwandte \textsc{Gauss} auf den Vergleich seiner Ergebnisse mit den Beobachtungen
bzw. auf die Bestimmung der Normalelemente, die in Verbindung mit den berechneten Störungen die
Beobachtungen möglichst gut darstellten. So lässt sich sehr gut erkennen,
wie genau die Beobachtungen mit der Theorie übereinstimmen. Die Abweichungen zwischen beiden
hat \textsc{Gauss} meist schon selbst veröffentlicht (s.~Band VI). Man findet sie außerdem bei den Vergleichen
mit den Oppositionen und in den oben abgedruckten Briefstellen. Zusätzlich zu den im Anschluss an die Bestimmung
der Störungen oben abgedruckten Vergleichen mit den Oppositionen findet sich im Nachlass
eine sehr große Menge solcher Vergleiche, auch für spätere Oppositionen, teils mit Berücksichtigung
der Saturnstörungen, teils ohne sie, teils mit dem alten, teils mit einem verbesserten Wert der Jupitermasse.
Am vollständigsten ist ein Vergleich mit den ersten 10 bzw. 12 Oppositionen (1802 bis 1817) im Handbuch
Bf erhalten. Er ist nicht abgedruckt, weil er methodisch nichts Neues bietet und seine
Ergebnisse aus dem auf S.~430 abgedruckten Brief an \textsc{Olbers} hervorgehen. Diese Vergleiche mit den Oppositionen
hat \textsc{Gauss} sowohl für seine speziellen Störungen durchgeführt, deren Berechnung er noch bis 1816
fortgesetzt zu haben scheint, als auch für die allgemeinen Störungen. Bei der Berechnung der Störungsbeträge aus
den Ausdrücken für die allgemeinen Störungen haben ihm \textsc{Westphal}, \textsc{Encke} und \textsc{Nicolai} wesentliche
Hilfe geleistet. Dazu wurden, solange die Tafeln (S.~572) noch nicht fertig waren, Gruppen von Störungsgleichungen
je nach ihrer Größe auf kleinen Zetteln zusammengefasst und »Manipel« genannt. Jeder Manipel
(mit Ausnahme des letzten) enthielt zehn Störungsglieder. So lauten z.~B. die ersten Manipel für die
Länge des Perihels wie folgt:

\needspace{19\baselineskip}
\begin{center}\textit{Länge des Perihels.}\end{center}
\begingroup\small
\setlength{\arraycolsep}{3pt}
\begin{center}\textit{Manipel 1.}\end{center}
\[
\begin{array}{ccc}
\left[\begin{gathered}\text{\itshape Nr. der Stö-}\\\text{\itshape rungsgleichung}\end{gathered}\right]&
\text{\itshape [Argument]}&
\left[\begin{gathered}\text{\itshape Log. des}\\\text{\itshape Koeffizienten}\end{gathered}\right]\\
\mathit{27}&\mathit{2}\BIntroJ-\BIntroP+\mathit{236^0\,43'\,49''}&\mathit{2{,}88715}\\
\mathit{66}&\mathit{5}\BIntroJ-\mathit{2}\BIntroP+\mathit{303\ \,26\ \,18}&\mathit{2{,}81492}\\
\mathit{40}&\mathit{3}\BIntroJ-\BIntroP+\mathit{94\ \,47\ \,46}&\mathit{2{,}72535}\\
\mathit{52}&\mathit{4}\BIntroJ-\mathit{2}\BIntroP+\mathit{288\ \,15\ \,4}&\mathit{2{,}16040}\\
\mathit{39}&\mathit{3}\BIntroJ-\mathit{2}\BIntroP+\mathit{260\ \,27\ \,22}&\mathit{1{,}95730}\\
\mathit{26}&\mathit{2}\BIntroJ-\mathit{2}\BIntroP+\mathit{342\ \,9\ \,24}&\mathit{1{,}94234}\\
\mathit{28}&\mathit{2}\BIntroJ\hspace{2em}+\mathit{101\ \,42\ \,47}&\mathit{1{,}91320}\\
\mathit{100}&\mathit{8}\BIntroJ-\mathit{3}\BIntroP+\mathit{168\ \,38\ \,33}&\mathit{1{,}90406}\\
\mathit{53}&\mathit{4}\BIntroJ-\BIntroP+\mathit{106\ \,12\ \,53}&\mathit{1{,}88715}\\
\mathit{14}&\BIntroJ\hspace{2em}+\mathit{291\ \,57\ \,59}&\mathit{1{,}85959}
\end{array}
\]
\begin{center}\textit{Manipel 2.}\end{center}
\[
\begin{array}{ccc}
\left[\begin{gathered}\text{\itshape Nr. der Stö-}\\\text{\itshape rungsgleichung}\end{gathered}\right]&
\text{\itshape [Argument]}&
\left[\begin{gathered}\text{\itshape Log. des}\\\text{\itshape Koeffizienten}\end{gathered}\right]\\
\mathit{78}&\mathit{6}\BIntroJ-\mathit{2}\BIntroP+\mathit{137^0\,59'\,40''}&\mathit{1{,}77929}\\
\mathit{13}&\BIntroJ-\BIntroP+\mathit{309\ \,34\ \,12}&\mathit{1{,}76702}\\
\mathit{1}&\hspace{1.5em}\BIntroP+\mathit{321\ \,28\ \,57}&\mathit{1{,}75857}\\
\mathit{89}&\mathit{7}\BIntroJ-\mathit{3}\BIntroP+\mathit{334\ \,36\ \,27}&\mathit{1{,}75507}\\
\mathit{38}&\mathit{3}\BIntroJ-\mathit{3}\BIntroP+\mathit{340\ \,39\ \,45}&\mathit{1{,}68732}\\
\mathit{77}&\mathit{6}\BIntroJ-\mathit{3}\BIntroP+\mathit{320\ \,40\ \,29}&\mathit{1{,}59337}\\
\mathit{65}&\mathit{5}\BIntroJ-\mathit{3}\BIntroP+\mathit{306\ \,34\ \,54}&\mathit{1{,}57175}\\
\mathit{122}&\mathit{10}\BIntroJ-\mathit{4}\BIntroP+\mathit{18\ \,48\ \,20}&\mathit{1{,}54904}\\
\mathit{41}&\mathit{3}\BIntroJ\hspace{2em}+\mathit{87\ \,0\ \,49}&\mathit{1{,}50939}\\
\mathit{12}&\BIntroJ-\mathit{2}\BIntroP+\mathit{242\ \,28\ \,11}.&\mathit{1{,}43467}
\end{array}
\]
\endgroup
\begin{center}\textit{usw.}\end{center}

\par\noindent Für jeden Manipel wurden die einzelnen Störungsglieder berechnet und dann die Beträge der
Manipel addiert, um die Gesamtstörungen jedes Elements zu bestimmen. Spezielle Störungen für spätere
Jahre hat auch \textsc{Encke} berechnet, s.~Berliner Astronomisches Jahrbuch für 1837—38.

\par\noindent\hspace*{\parindent}Die erste Untersuchung zur Jupitermasse ist auf S.~561 abgedruckt. Der Verlauf von \textsc{Gauss}’ Arbeiten
daran ergibt sich aus dem Briefwechsel auf S.~424—429. Auffallend ist, wie genau \textsc{Gauss} den Wert
dieser Größe zuletzt im Brief an \textsc{Olbers} vom 24.~Juli 1816 mit \(\frac{1}{1030}\) angibt.

\par\noindent\hspace*{\parindent}Die Briefe zu den auf S.~572 abgedruckten Tafeln der Störungen durch Jupiter stammen aus der Zeit von
\SourceLeaf{609}Februar 1817 bis März 1818. Der erste Plan für eine solche Tafel wird bereits im Brief an \textsc{Olbers} vom
8.~April 1813 erwähnt, in welchem \textsc{Gauss} auch die Absicht ausspricht, die Störungen durch Saturn und
Mars berechnen zu lassen.

\par\noindent\hspace*{\parindent}Hinweise auf die Ausarbeitung seiner Methode, als deren Beginn wir die »Exposition d’une nouvelle méthode
etc.« ansehen müssen, finden sich in den Briefen vom 12.~Dezember 1813 und vom 25.~September
1814. Dass sie in der Absicht begonnen wurde, sich um den von der französischen Akademie ausgesetzten
Preis zu bewerben, geht schon daraus hervor, dass sie auf Französisch abgefasst wurde. Wie sich
dann die Hoffnungen der astronomischen Welt, der Preis möchte \textsc{Gauss} zufallen, durch lange Jahre hinzogen
und schließlich doch enttäuscht wurden, zeigt der Briefwechsel deutlich genug.
\endgroup
\endgroup
